% Part: incompleteness
% Chapter: representability-in-q
% Section: representable-comp

\documentclass[../../../include/open-logic-section]{subfiles}

\begin{document}

\olfileid{inc}{req}{rpc}
\olsection{$\Th{Q}$ मध्ये निरूपणीय फलने संगणनक्षम असतात}

$\Th{Q}$ मध्ये निरूपणीय असलेले प्रत्येक फलन संगणनक्षम आहे,
हे आपण सिद्ध करू. त्याआधी $\Th{Q}$ मध्ये निरूपणीय फलनांविषयी
एक साहाय्यक विधान सिद्ध करावे लागेल.

\begin{lem}\ollabel{lem:rep-q}
  $f(x_0,
  \dots, x_k)$ हे $\Th{Q}$ मध्ये निरूपणीय असेल, तर
  असे !!a{formula}~$!A_f(x_0, \dots, x_k, y)$ असते की
  \[
  \Th{Q} \Proves !A_f(\num{n_0}, \dots, \num{n_k}, \num{m})
  \quad\text{iff}\quad m = f(n_0, \dots, n_k).
  \]
\end{lem}

\begin{proof}
  ``तर'' हा भाग
\olref[int]{defn:representable-fn}\olref[int]{defn:rep:a} मधून मिळतो.
``तरच'' हा भाग पुढीलप्रमाणे दिसतो: समजा $\Th{Q} \Proves !A_f(\num{n_0},
\dots, \num{n_k}, \num{m})$, पण $m \neq f(n_0, \dots, n_k)$.
$l = f(n_0, \dots, n_k)$ असे धरू.
\olref[int]{defn:representable-fn}\olref[int]{defn:rep:a} नुसार $\Th{Q}
\Proves !A_f(\num{n_0}, \dots, \num{n_k}, \num{l})$.
\olref[int]{defn:representable-fn}\olref[int]{defn:rep:b} नुसार
$\lforall[y][(!A_f(\num{n_0}, \dots, \num{n_k}, y) \lif \num{l} =
y)]$. तर्कशास्त्र वापरून आणि $\Th{Q} \Proves
!A_f(\num{n_0}, \dots, \num{n_k}, \num{m})$ या गृहीतकावरून
$\Th{Q} \Proves \eq[\num{l}][\num{m}]$ मिळते. दुसरीकडे,
\olref[bre]{lem:q-proves-neq} नुसार $\Th{Q} \Proves
\eq/[\num{l}][\num{m}]$. म्हणून $\Th{Q}$ विसंगत ठरेल.
पण ते अशक्य आहे: मानक प्रतिमान $\Th{Q}$ ची पूर्ती करते
(पाहा \olref[int][def]{def:standard-model}), म्हणजे
$\Sat{N}{\Th{Q}}$; आणि निर्दोषता प्रमेयानुसार प्रत्येक
पूर्ततायोग्य उपपत्ती सुसंगत असते
(\tagrefs{prfAX/{fol:axd:sou:cor:consistency-soundness},
prfSC/{fol:seq:sou:cor:consistency-soundness},
prfND/{fol:ntd:sou:cor:consistency-soundness},
prfTab/{fol:tab:sou:cor:consistency-soundness}}).
\end{proof}

\begin{lem}
$\Th{Q}$ मध्ये निरूपणीय असलेले प्रत्येक फलन संगणनक्षम आहे.
\end{lem}

\begin{proof}
हे खरे असण्यामागील अंतर्ज्ञानी कल्पना आधी पाहू. $f$ चे मूल्य
काढण्यासाठी पुढील पद्धत वापरू. अंकगणिताच्या भाषेतील सर्व संभाव्य
!!{derivation}s~$\delta$ यांची यादी करायची; हे यांत्रिकरीत्या
करता येते. प्रत्येकासाठी, ती $!A_f(\num{n_0}, \dots,
\num{n_k}, \num{m})$ या रूपाच्या !!a{formula} ची
!!a{derivation} आहे का ते तपासायचे (हे $f$ चे $\Th{Q}$
मधील निरूपण करणारे !!{formula} आहे; पाहा \olref{lem:rep-q}).
असेल, तर \olref{lem:rep-q} नुसार $m = f(n_0, \dots, n_k)$;
म्हणजे $f$ चे मूल्य सापडले. हा शोध संपतो, कारण
$\Th{Q} \Proves !A_f(\num{n_0}, \dots, \num{n_k},
\num{f(n_0, \dots, n_k)})$; त्यामुळे योग्य प्रकारची
$\delta$ अखेरीस सापडते.

ही मांडणी अजून पूर्णपणे नेमकी नाही: आपली पद्धत फक्त संख्यांवर
न चालता !!{derivation}s आणि !!{formula}s वर चालते, आणि
``सर्व संभाव्य !!{derivation}s ची यादी'' यांत्रिकरीत्या कशी
करता येते हे आपण अजून स्पष्ट केलेले नाही. पण आपण पाहिलेच आहे
की पदे, !!{formula}s आणि !!{derivation}s यांना गोडेल संख्यांनी
संकेतबद्ध करता येते. संगणनाचे नेमके प्रतिमान म्हणून सामान्य
पुनरावर्ती फलनेही आपण मांडली आहेत. तसेच
$\Prf[\Th{Q}](d,y)$ हा संबंध (आदिम) पुनरावर्ती आहे:
नैसर्गिक निगमनातील थेट निष्पत्ती-संकेतांकाचा अर्थ घेतल्यास,
$d$ ही $\Th{Q}$ च्या स्वयंसिद्धकांपासून $y$ ही गोडेल संख्या
असलेल्या !!{formula} च्या !!a{derivation} ची गोडेल संख्या
असेल तेव्हाच हा संबंध खरा असतो. आपल्याला आणखी लागणारी आदिम
पुनरावर्ती फलने म्हणजे $\fn{num}$
(\olref[art][trm]{prop:num-primrec}) आणि $\fn{Subst}$
(\olref[art][sub]{prop:subst-primrec}). यांपासून लघुतमीकरणाने
$f$ परिभाषित करता येते; म्हणून $f$ पुनरावर्ती आहे.

प्रथम पुढील व्याख्या करा:
\begin{multline*}
  A(n_0, \dots, n_k, m) = \\
  \fn{Subst}(\fn{Subst}(\dots\fn{Subst}(\Gn{!A_f}, \fn{num}(n_0), \Gn{x_0}),\\
  \dots),
  \fn{num}(n_k),  \Gn{x_k}), \fn{num}(m), \Gn{y})
\end{multline*}
हे गुंतागुंतीचे दिसते, पण हे फक्त $A(n_0, \dots, n_k,
m) = \Gn{!A_f(\num{n_0}, \dots, \num{n_k}, \num{m})}$
हे फलन आहे.

आता $R(n_0, \dots, n_k, s)$ हा संबंध पाहा. $(s)_0$ ही
$!A_f(\num{n_0}, \dots, \num{n_k}, \num{(s)_1})$ च्या
$\Th{Q}$ पासूनच्या !!a{derivation} ची गोडेल संख्या असेल
तेव्हा हा संबंध खरा असतो:
\[
R(n_0, \dots, n_k, s) \quad\text{iff}\quad \Prf[\Th{Q}]((s)_0, A(n_0,
\dots, n_k, (s)_1))
\]
$R(n_0, \dots, n_k, s)$ खरा असणारा एखादा $s$ मिळाला,
तर $(s)_0$ आणि $(s)_1$ ही संख्यांची जोडी मिळाली आहे:
$(s)_0$ ही $!A_f(\num{n_0}, \dots,
\num{n_k}, \num{(s)_1})$ च्या !!a{derivation} ची गोडेल संख्या
असते. म्हणून $s$ शोधणे म्हणजे अनौपचारिक पुराव्यातील $d$ आणि
$m$ ही जोडी शोधण्यासारखे आहे. असा $s$ ``शोधणारे'' संगणनक्षम
फलन नियमित लघुतमीकरणाने परिभाषित करता येते. लक्षात घ्या,
$R$ साठी हा शोध नियमित आहे: प्रत्येक $n_0$, \dots, $n_k$
साठी $\Th{Q} \Proves !A_f(\num{n_0}, \dots,
\num{n_k}, \num{f(n_0, \dots, n_k)})$ ची !!a{derivation}~$\delta$
असते. म्हणून $s = \tuple{\Gn{\delta}, f(n_0, \dots, n_k)}$
घेतल्यास $R(n_0, \dots, n_k, s)$ खरा असतो. त्यामुळे $f$
पुढीलप्रमाणे लिहिता येते:
\[
f(n_0,\dots,n_{k}) = (\umin{s}{R(n_0, \dots, n_k, s)})_1.
\]
\end{proof}

\end{document}
